\section{İlgili Çalışmalar}
\label{sec:related_work}

\emph{RIS faz optimizasyonu.} Ortak ön kodlama ve faz tasarımı; dönüşümlü optimizasyon ve yarı tanımlı gevşetme~\cite{Wu2019TWC}, faza bağımlı pratik genlik modelleri~\cite{Abeywickrama2020TCOM}, blok koordinat inişi~\cite{Pan2020TWC}, kesirli programlama~\cite{Guo2020TWC}, sabit nokta ve manifold yöntemleri~\cite{Yu2019ICCC} ve~\cite{Shi2011TSP} çalışmasındaki algoritmanın üzerine kurulan ağırlıklı en küçük ortalama karesel hata (WMMSE) yöntemleri~\cite{Choi2024TCOM} ile ele alınmıştır. Bu çalışmalar ulaşılan amaç fonksiyonu değerini raporlar; koşma süresi genellikle bir karmaşıklık mertebesi olarak verilir ve ölçüldüğü durumlarda da sabit kanal gerçekleşmeleri üzerinde ölçülür~\cite{Abeywickrama2020TCOM,Yu2019ICCC,Shi2011TSP,Cao2022TWC}; eş evrelilik süresiyle karşılaştırılmaz.

\emph{Öğrenme.} Öğrenilmiş politikalar arasında ortak ön kodlama ve faz tasarımı için derin pekiştirmeli öğrenme~\cite{Huang2020JSAC} ve yansıtmayı doğrudan pilotlardan öğrenme~\cite{Jiang2021JSAC} yer alır. Algoritma açma (unrolling), bir optimize edicinin yinelemelerini ağ katmanlarına dönüştürür~\cite{Monga2021SPM}; bu yaklaşım RIS hüzme oluşturmaya~\cite{Liu2022TWC}, MIMO sezimine~\cite{Samuel2019TSP} ve WMMSE ön kodlamaya~\cite{Hu2021TWC,Pellaco2022OJCOMS} uygulanmıştır ve eğitilmiş ağlar WMMSE'yi onun koşma süresinin küçük bir kesrinde yaklaşık olarak gerçekleyebilir~\cite{Sun2018TSP}. Her fazın sinüs ve kosinüsünü öngörmek, doğrudan faz regresyonunun $\pm\pi$ süreksizliğini ortadan kaldırır~\cite{SavasciSen2026AISY}. Öğrenilen yöntemler genellikle hızla gerekçelendirilir, ancak bu hız nadiren kanalın ne kadar hızlı değiştiğiyle ilişkilendirilir; her ikisi tek bir makinede zamanlandığında bile~\cite{Samuel2019TSP,Sun2018TSP,Hu2021TWC,Liu2022TWC}, klasik taban çizgisi aynı gecikmeye kadar değil, yakınsamaya ya da sabit bir yineleme sayısına kadar koşturulur.

\emph{Hareketlilik ve güncelliğini yitirmiş kanallar.} İki aşamalı kestirim kullanan, araca monte edilmiş kırıcı bir yüzey, hızlı sönümlemeyi yavaş sönümlemeye dönüştürür~\cite{Huang2022TWC}. RIS için kanal eskimesi~\cite{Papazafeiropoulos2023TVT} ve güncelliğini yitirmiş kanal bilgisi~\cite{Zhang2022TCOM} analiz edilmiştir; öngörücü hüzme oluşturma, hüzmeleri hızlı insansız hava araçlarının öngörülen konumlarına yöneltir~\cite{Liu2021WCL} ve etkin algılama hareketli kullanıcıları izler~\cite{Han2025TWC}. Eskime, büyük MIMO'da klasik bir konudur~\cite{Truong2013JCN,Chopra2018TWC}; yayın kanallarında sınırlı kanal geri beslemesi de öyledir~\cite{Jindal2006IT}. İki zaman ölçekli tasarımlar fazları kanal istatistiklerinden, ön kodlayıcıyı ise anlık kanallardan günceller~\cite{Zhao2021TWC,Cao2022TWC}. Pratik yüzeylerin birkaç bitlik faz çözünürlüğü vardır~\cite{Wu2020TCOM,Zhao2024OJCOMS}; Bölüm~\ref{sec:res_robust}, sürekli faz sonuçlarımıza 1--3 bitlik fazları ekler. Bu çalışmalarda kanal edinimi ile kullanımı arasındaki gecikme, modellendiği yerde sabittir; hesaplama bir neden olarak anıldığı yerde bile~\cite{Truong2013JCN} böyledir; optimize edicinin kendi koşma süresi en fazla bir uygulanabilirlik kontrolü ya da sabit bir ek yüktür~\cite{Ryu2026arXiv,Liu2021WCL} ve eskimeye dahil edilmez.

\emph{Boşluk.} Bildiğimiz kadarıyla hiçbir RIS çalışması her yöntemi kendi ölçülmüş tamamlanma anındaki kanal üzerinde değerlendirmemekte, bunun optimal bir yineleme sayısı ortaya çıkardığını göstermemekte ya da öğrenilen ve klasik optimize edicileri hem eşit adımda hem de eşit ölçülmüş sürede karşılaştırmamaktadır.
